%% Annexe B : Référentiel des outils
\chapter{Référentiel des Outils}
\label{ann:outils}

Les outils cités sont des exemples. Le cours insiste sur la méthode~: l'outil
sert l'enquête, il ne la remplace pas. Tout usage se fait dans un cadre
autorisé, sur des copies.

\begin{tableaucours}{L{4.2cm}L{5.4cm}L{2.8cm}}{%
    \tableauHeader{Catégorie d'outil} &
    \tableauHeader{Usage} &
    \tableauHeader{Chapitre, lab}
}
Bloqueur en écriture & Lire un support sans l'altérer & Ch.~\ref{chap:preserver} \\
Duplication et hachage & Image bit à bit, condensat (SHA-256) & Ch.~\ref{chap:duplication}, FORN-LAB-02 \\
Examen de disque & Systèmes de fichiers, artefacts & Ch.~\ref{chap:disque}, FORN-LAB-03 \\
Récupération par signatures & Fichiers supprimés, fragments & Ch.~\ref{chap:artefacts}, FORN-LAB-04 \\
Analyse mémoire & Processus, connexions, secrets & Ch.~\ref{chap:memoire}, FORN-LAB-05 \\
Analyse de flux & Journaux de flux exportés & Ch.~\ref{chap:reseau}, FORN-LAB-06 \\
Extraction mobile & Messages, applications, géolocalisation & Ch.~\ref{chap:mobile}, FORN-LAB-07 \\
Tableur & Chronologie, faisceau, qualité & Ch.~\ref{chap:reconstitution}, FORN-LAB-09 \\
\end{tableaucours}

\begin{boxalert}
\textbf{Rappel.} Aucun outil ne s'applique à un support réel sans titre.
L'analyse se fait sur des copies vérifiées par condensat, dans le périmètre de
la réquisition.
\end{boxalert}
